\section{Implementation and On-Glasses Deployment}\label{sec:impl}

\subsection{Training}
We train with AdamW~\cite{loshchilov2019decoupled} (weight decay $10^{-4}$, batch size 32) and a one-cycle schedule~\cite{smith2019super} (peak learning rate $3\times10^{-3}$ for the heads and $3\times10^{-4}$ for the backbone, 10\% warm-up, then cosine decay) for up to 180 epochs, stopping early after 25 epochs without a lower validation loss, with gradient clipping at 1.0, mixed precision and seed 42. Training took 1\,h\,40\,min on one NVIDIA Tesla T4 (47\,s per epoch). Early stopping ended it after 125 epochs, before the one-cycle schedule had completed; the reported checkpoint is from epoch 100. Every run in this paper uses the same stopping rule, but runs stop at different epochs, which we take into account when comparing them (Section~\ref{sec:ablations}).

Besides the distractor pasting of Section~\ref{sec:ownership}, we use horizontal flips (but no vertical flips, because a wearer's hand entering from the top would contradict the ownership prior), shift, scale and rotation, perspective warps, brightness and contrast changes, Gaussian blur, and mixup~\cite{zhang2018mixup} between images of the same scene. All augmentation is applied on the fly, and Appendix~\ref{app:aug} lists all parameters. The model is exported to ONNX (opset 17, which bilinear grid sampling requires) and to FP16; the FP16 weights take 1.6\,MB.

\subsection{Deployment on Smart Glasses}\label{sec:device}
We deploy the 160\,px model on an engineering unit of a pair of smart glasses: a pre-production test device, not a retail product. It has a smart-glasses system-on-chip paired with a low-power microcontroller; a single 12\,MP ultra-wide camera (105$^\circ$) in portrait orientation, with a dedicated low-power vision NPU; Wi-Fi~6 and Bluetooth~5.4 to a companion phone; and a 160\,mAh battery, in a frame weighing under 45\,g. Because it is an engineering device, its power management is not tuned for production, and its pre-production software stack is not public. We therefore treat the unit as a measurement platform: we report end-to-end device measurements rather than runtime-specific details, and we compare power between operating modes on the same unit rather than reporting battery life.

We compare two ways of running pointing localisation. \emph{On the glasses}, SpatialNet runs continuously on frames from the glasses' camera. In \emph{offloading}, the glasses only capture frames and transmit them to the companion phone at 0.5 frames per second, the alternative in which localisation would run on the phone or in the cloud. We measure the inference latency per frame; the peak additional RAM used during inference, over the idle application with the model already loaded (the 1.6\,MB of FP16 weights are stored separately); and the battery drain over one minute in each operating mode, starting from the same charge level.

The camera delivers portrait frames, whereas most EgoPoint images are landscape. The model squashes both to a square, so the deployed input has a different aspect ratio from most training images. Section~\ref{sec:realframes} checks the model's accuracy on frames from this camera.
